\]

\begin{idea}
Saint-Venant justifica una gran parte de la Resistencia de Materiales: en vez de describir cada contacto real, tornillo o unión en detalle, trabajamos con sus resultantes.
\end{idea}

\safesection{Existencia y unicidad}
Bajo unas condiciones de contorno adecuadas y una ley constitutiva elástica estable, el problema lineal bien planteado posee una solución única.

En la práctica esto exige eliminar los movimientos de sólido rígido.

Por ejemplo, una viga simplemente apoyada está correctamente vinculada porque no puede trasladarse libremente como cuerpo rígido.

% ============================================================
\safechapter{Tensión}
% ============================================================

\safesection{Por qué aparece el concepto de tensión}
Supongamos un cuerpo en equilibrio.

Realizamos un corte imaginario de normal $\vect{n}$.

La parte eliminada ejercía sobre la parte restante fuerzas distribuidas a través de la superficie de corte.

La intensidad local de esas fuerzas se caracteriza mediante el \textbf{vector tensión}:
\[
\vect{t}^{(\vect n)}
=
\lim_{\Delta A\rightarrow 0}
\frac{\Delta \vect F}{\Delta A}.
\]

Sus unidades son fuerza dividida por área:
\[
[\text{tensión}]=\frac{\mathrm{N}}{\mathrm{m^2}}=\mathrm{Pa}.
\]

Una unidad extremadamente cómoda en estructuras es:
\[
1\ \mathrm{MPa}=1\ \mathrm{N/mm^2}.
\]

\safesubsection{Nota sobre una conversión habitual}
En algunos apuntes clásicos aparece aproximadamente
\[
1\ \mathrm{MPa}\simeq 10\ \mathrm{kp/cm^2}.
\]
El valor más preciso es
\[
1\ \mathrm{MPa}\simeq 10.197\ \mathrm{kp/cm^2}.
\]
Para cálculo moderno conviene trabajar directamente en SI.

\safesection{Tensión normal y tangencial}
\index{tensión!vector tensión}El vector tensión puede descomponerse en una componente normal y otra tangencial.

La componente normal es
\[
\sigma_n=\vect{t}^{(\vect n)}\cdot\vect n
\]
y la parte tangencial:
\[
\vect{\tau}=\vect{t}^{(\vect n)}-\sigma_n\vect n.
\]

Así distinguimos:
\[
\text{tensión normal }\sigma
\qquad\text{y}\qquad
\text{tensión tangencial }\tau.
\]

La tensión normal está asociada principalmente a tracción/compresión y flexión. La tensión tangencial aparece de forma importante en cortadura y torsión.

\safesection{Tensor de tensiones}
La tensión depende de la orientación del plano de corte.

El teorema de Cauchy conduce a
\[
\vect t^{(\vect n)}=\mat{\sigma}\vect n
\]
donde
\[
\mat{\sigma}=
\begin{bmatrix}
\sigma_{xx}&\tau_{xy}&\tau_{xz}\\
\tau_{yx}&\sigma_{yy}&\tau_{yz}\\
\tau_{zx}&\tau_{zy}&\sigma_{zz}
\end{bmatrix}.
\]

En ausencia de pares internos clásicos, el equilibrio angular implica
\[
\tau_{xy}=\tau_{yx},\qquad
\tau_{xz}=\tau_{zx},\qquad
\tau_{yz}=\tau_{zy}.
\]

\safesection{Tensión nominal y tensión verdadera}
En un ensayo uniaxial:
\[
\sigma_{\mathrm{nom}}=\frac{F}{A_0}
\]
utiliza el área inicial.

La tensión verdadera utiliza el área instantánea:
\[
\sigma_{\mathrm{real}}=\frac{F}{A}.
\]

Cuando la deformación es pequeña, $A\simeq A_0$ y ambas son prácticamente iguales.

% ============================================================
\safechapter{Deformación}
% ============================================================

\safesection{Desplazamiento no es deformación}
Un cuerpo puede desplazarse sin deformarse.

Por ejemplo, una traslación rígida
